\documentclass[11pt,a4paper]{article}
\usepackage[margin=26mm]{geometry}
\usepackage{amsmath,amssymb,amsthm,mathtools}
\usepackage{lmodern,microtype,booktabs,array}
\usepackage[colorlinks=true,linkcolor=blue,citecolor=blue,urlcolor=blue]{hyperref}
\hypersetup{pdftitle={Local zero-free shapes from reflected logarithmic-derivative kernels},pdfauthor={Li Xiang}}
\newtheorem{theorem}{Theorem}[section]
\newtheorem{proposition}[theorem]{Proposition}
\newtheorem{lemma}[theorem]{Lemma}
\newtheorem{corollary}[theorem]{Corollary}
\theoremstyle{definition}
\newtheorem{definition}[theorem]{Definition}
\newtheorem{remark}[theorem]{Remark}
\newcommand{\Rea}{\operatorname{Re}}
\newcommand{\Ima}{\operatorname{Im}}
\newcommand{\scen}{\sigma_c}
\newcommand{\Tcen}{T_0}
\newcommand{\src}[1]{\nolinkurl{#1}}
\makeatletter
\g@addto@macro\UrlBreaks{\do\a\do\b\do\c\do\d\do\e\do\f\do\g\do\h\do\i\do\j\do\k\do\l\do\m\do\n\do\o\do\p\do\q\do\r\do\s\do\t\do\u\do\v\do\w\do\x\do\y\do\z\do\0\do\1\do\2\do\3\do\4\do\5\do\6\do\7\do\8\do\9}
\makeatother
\title{Local zero-free shapes from reflected logarithmic-derivative kernels}
\author{Li Xiang}
\date{October 11, 2026}

\begin{document}
\maketitle

\begin{abstract}
We construct local additions to a zero-free half-plane for the Riemann
zeta function, conditional on the global bound $\beta\le 7/8$ for its
nontrivial zeros. Positivity of the Hadamard real-part expansion to the
right of this bound gives a single-zero disk. Keeping the same-height
reflected zero gives a larger exclusion region described by a quartic
inequality. An external Arb computation at
$219/250+4\cdot10^{12}i$ supplies the upper bound $29/2$ for the real
logarithmic derivative of the completed function. Exact rational corner
values and concavity then exclude every point with
$\Rea s\ge4/5$ and $|\Ima s-4\cdot10^{12}|\le1/200$.
The real algebra and conditional exclusion implications are formalized
in Lean. The global zero bound, the analytic kernel identity, and the
Arb evaluation remain separate inputs. We also give a finite nonnegative
multi-anchor certificate, with its pointwise limitation made explicit.
The result is a local shape extension; no new uniform strip or
end-to-end formalized zeta theorem is claimed.
\end{abstract}

\section{Statement and scope}
\label{sec:scope}

Let
\begin{equation}\label{eq:xi}
 \xi(s)=\frac12 s(s-1)\pi^{-s/2}\Gamma(s/2)\zeta(s).
\end{equation}
This is the entire completed Riemann function. Its zeros are precisely
the nontrivial zeros of $\zeta$, with the same multiplicities. Throughout
the paper, a sum over zeros includes all ordinates of both signs and
includes multiplicities.

Our global analytic hypothesis is
\begin{equation}\tag{$H_{7/8}$}\label{eq:H}
 \zeta(\beta+i\gamma)=0,\quad 0<\beta<1
 \quad\Longrightarrow\quad \beta\le\frac78.
\end{equation}
\addtocounter{equation}{1}
The seven-eighths input and the stronger cubic boundary studied in
\cite{Cubic} are used only with their stated proof dependencies. In
particular, this paper does not supply the imported arithmetic probe
obligation of that work. The local construction below needs only
\eqref{eq:H}; it is independent of the small improvement from $7/8$ to
the cubic boundary.

Fix the exact constants
\begin{equation}\label{eq:constants}
 \scen=\frac{219}{250},\qquad
 \Tcen=4000000000000,\qquad E=\frac{29}{2},\qquad
 K=\Rea\frac{\xi'}{\xi}(\scen+i\Tcen).
\end{equation}
Notice that $\scen-7/8=1/1000>0$. Thus \eqref{eq:H}
ensures that $\xi(\scen+i\Tcen)\ne0$ and that $K$ is well defined.

\begin{theorem}[Conditional local rectangle]\label{thm:rectangle}
Assume \eqref{eq:H} and $K\le E$, with the constants in
\eqref{eq:constants}. Then
\begin{equation}\label{eq:rectangle}
 \zeta(s)\ne0\qquad\text{if}\qquad
 \Rea s\ge\frac45,\qquad
 \left|\Ima s-4000000000000\right|\le\frac1{200}.
\end{equation}
The same conclusion holds with the central ordinate replaced by
$-4000000000000$. All displayed rectangle edges are included.
\end{theorem}

Section~\ref{sec:numeric} records a rigorous external ball-arithmetic
calculation of $K<E$, conditional only on accepting the correctness of
the cited numerical library and the executable evaluation. Thus the
numerical instance of Theorem~\ref{thm:rectangle} still depends on
\eqref{eq:H}. Theorem~\ref{thm:quartic} gives the curved region of which
\eqref{eq:rectangle} is a particularly simple rational subset.

The choice of $\Tcen$ is an illustration at a high ordinate. It exceeds
the height $3000175332800$ covered by the computation in
Platt--Trudgian \cite{PlattTrudgian}; we make no assertion about the
largest verified height at the date of this paper. Our calculation does
not enumerate zeros or verify the Riemann hypothesis up to $\Tcen$.

\section{Positive kernels from the completed function}
\label{sec:kernel}

We first recall the analytic identity underlying the construction.
Related real-part positivity questions are discussed in
Gold\v{s}tein--Grigutis \cite{GoldsteinGrigutis}. The identity used here
comes from the classical Hadamard product; the local geometric
consequences will be derived explicitly.

\begin{lemma}[Real-part Hadamard identity]\label{lem:hadamard}
For any $s=\sigma+it$ with $\xi(s)\ne0$,
\begin{equation}\label{eq:hadamard}
 \Rea\frac{\xi'}{\xi}(s)
 =\sum_{\rho=\beta+i\gamma}
 \frac{\sigma-\beta}{(\sigma-\beta)^2+(t-\gamma)^2}.
\end{equation}
The real series on the right is absolutely convergent.
\end{lemma}

\begin{proof}
The order-one Hadamard product gives
\[
 \frac{\xi'}{\xi}(s)
 =b_0+\sum_\rho\left(\frac1{s-\rho}+\frac1\rho\right)
\]
with the canonical summation convention. For fixed $s$ and large
$|\gamma|$, the real part of each of $1/(s-\rho)$ and $1/\rho$
is $O_s(\gamma^{-2})$, since $0<\beta<1$. The classical zero count
$N(T)=O(T\log T)$ therefore gives absolute convergence of the separate
real-part series. Taking real parts leaves a possible real constant
\[
 C=\Rea b_0+\sum_\rho\Rea\frac1\rho.
\]
This constant is zero. Indeed, the functional equation
$\xi(s)=\xi(1-s)$ gives $\xi'(1/2)=0$, while $\xi(1/2)\ne0$.
The latter is also immediate from the alternating-series expression
for $\zeta(1/2)$, which is negative. At $s=1/2$, the terms
$\Rea(1/(1/2-\rho))$ cancel under
$\rho\mapsto1-\overline\rho$; critical-line terms are zero.
Absolute convergence justifies this pairing. Hence $C=0$, proving
\eqref{eq:hadamard}. This argument also shows that no additional real
constant or factor of two belongs in \eqref{eq:hadamard}.
\end{proof}

More generally, assume that every nontrivial zero satisfies
$\beta\le B$, where $1/2\le B<1$, and choose an anchor
$s_c=\sigma_c+iT$ with $\sigma_c>B$. Every summand in
\eqref{eq:hadamard} is then strictly positive. For an individual zero,
\begin{equation}\label{eq:single-major}
 \frac{\sigma_c-\beta}{(\sigma_c-\beta)^2+(T-\gamma)^2}
 \le \Rea\frac{\xi'}{\xi}(s_c).
\end{equation}
This use of positivity needs no truncation or estimate of the remaining
zero sum.

\begin{proposition}[Single-zero disk]\label{prop:disk}
Under the preceding global bound, suppose
$\Rea(\xi'/\xi)(\sigma_c+iT)\le E$ for a real $E>0$. Then every
point of the open disk
\begin{equation}\label{eq:disk}
 \left(\sigma-\left(\sigma_c-\frac1{2E}\right)\right)^2
 +(t-T)^2<\frac1{4E^2}
\end{equation}
is not a nontrivial zero of $\zeta$.
\end{proposition}

\begin{proof}
Set $a=\sigma_c-\sigma$ and $u=t-T$. The strict disk inequality is
$a^2+u^2<a/E$. It implies $a>0$ and $a^2+u^2>0$; in particular
the anchor, where $a=u=0$, is not inside the disk. Dividing gives
$a/(a^2+u^2)>E$, contrary to \eqref{eq:single-major} for a zero.
\end{proof}

For \eqref{eq:constants}, this disk has radius $1/29$, real center
$6101/7250$, and leftmost real coordinate $5851/7250$
($0.8070344827\ldots$). The reflected-pair construction below reaches
further left.

\section{The reflected-pair quartic region}
\label{sec:quartic}

\begin{lemma}[Two distinct same-height zeros]\label{lem:pair}
Let $\rho=\beta+i\gamma$ be a nontrivial zero with $\beta>1/2$.
Then $1-\overline\rho=1-\beta+i\gamma$ is a distinct zero of the
same multiplicity. Under the global bound $\beta\le B<\sigma_c$,
\begin{equation}\label{eq:pair-major}
 \Phi_{\sigma_c,T}(\beta,\gamma)
 :=\frac{\sigma_c-\beta}{(\sigma_c-\beta)^2+(\gamma-T)^2}
 +\frac{\sigma_c+\beta-1}{(\sigma_c+\beta-1)^2+(\gamma-T)^2}
 \le\Rea\frac{\xi'}{\xi}(\sigma_c+iT).
\end{equation}
\end{lemma}

\begin{proof}
Complex conjugation and the functional equation of $\xi$ preserve zero
multiplicities. Their composition gives the stated same-height
reflection. Distinctness follows from $\beta>1/2$. Both terms in
\eqref{eq:pair-major} occur in the positive sum
\eqref{eq:hadamard}. Even when a zero has multiplicity greater than
one, taking one copy of each gives the stated inequality; taking all
copies could only strengthen it. Critical-line zeros would coincide
with their reflected zeros and cannot be counted twice by this argument.
\end{proof}

For $1/2<\sigma\le B$ put
\begin{equation}\label{eq:abpz}
 a=\sigma_c-\sigma>0,\quad b=\sigma_c+\sigma-1>0,\quad
 c=a+b=2\sigma_c-1,\quad p=ab>0,\quad z=(t-T)^2\ge0.
\end{equation}
The sum of the two kernels simplifies to
\begin{equation}\label{eq:pair-identity}
 \frac a{a^2+z}+\frac b{b^2+z}
 =\frac{c(p+z)}{(p-z)^2+c^2z}.
\end{equation}
In fact, the common denominator is
\[
 (a^2+z)(b^2+z)=(p-z)^2+c^2z>0,
\]
and the numerator is $(a+b)(ab+z)$. The strict positivity of
the denominator also follows directly from $a,b>0$ and $z\ge0$.

\begin{theorem}[Quartic shape certificate]\label{thm:quartic}
Assume that all nontrivial zeros satisfy $\beta\le B$, where
$1/2\le B<1$ and $\sigma_c>B$, and that
$\Rea(\xi'/\xi)(\sigma_c+iT)\le E$ for $E>0$. Define
\begin{equation}\label{eq:surplus}
 S_E(p,z)=c(p+z)-E\bigl((p-z)^2+c^2z\bigr),
 \qquad c=2\sigma_c-1.
\end{equation}
Then the region
\begin{equation}\label{eq:shape-domain}
 \frac12<\sigma\le B,\qquad
 S_E\bigl((\sigma_c-\sigma)(\sigma_c+\sigma-1),(t-T)^2\bigr)>0
\end{equation}
contains no zeros of $\zeta$. Its union with $\{\Rea s>B\}$
contains no nontrivial zeros. The latter half-plane contains the pole
$s=1$, so it is a zero-free meromorphic domain rather than a domain
on which $\zeta$ is everywhere finite.
\end{theorem}

\begin{proof}
If a zero lay in \eqref{eq:shape-domain}, then
\eqref{eq:pair-identity}, the positive denominator, and $S_E>0$
would imply $\Phi_{\sigma_c,T}(\sigma,t)>E$.
This contradicts Lemma~\ref{lem:pair} and the assumed upper bound.
The half-plane assertion follows from the global bound.
\end{proof}

The polynomial in \eqref{eq:shape-domain} has total degree four in
$\sigma$ and $t-T$. This degree arises from two positive kernel
terms; it is not a fourth-moment estimate for $\zeta$ or for a
detector trace.

One can describe its horizontal slices without approximating a
boundary curve. For fixed $z\ge0$,
\[
 S_E(p,z)=-Ep^2+(c+2Ez)p+cz-Ez^2-Ec^2z.
\]
The discriminant is
\begin{equation}\label{eq:disc}
 D(z)=c^2+4Ecz(2-Ec).
\end{equation}
If $D(z)>0$, the condition is precisely
\begin{equation}\label{eq:root-domain}
 p_-(z)<p<p_+(z),\qquad
 p_\pm(z)=\frac{c+2Ez\pm\sqrt{D(z)}}{2E}.
\end{equation}
If $D(z)\le0$, no strict positive slice exists. These statements are
only to be used together with $1/2<\sigma\le B$ and $p>0$.
At $z=0$, \eqref{eq:root-domain} reduces to $0<p<c/E$.
For the constants \eqref{eq:constants}, the corresponding left edge
at the central ordinate is
\begin{equation}\label{eq:center-edge}
 \frac12+\sqrt{\left(\frac{219}{250}-\frac12\right)^2
                       -\frac{94/125}{29/2}}
 =0.7991887882\ldots.
\end{equation}
The exact radical, rather than this illustrative decimal, specifies
the edge. Strict inequality excludes points to its right within the
quartic domain.

\begin{remark}[Do not fill unproved gaps]\label{rem:gap}
A quartic component can be separated from the old half-plane at some
heights. The certified region is the union actually stated in
Theorem~\ref{thm:quartic}; it is not automatically the entire
half-plane to the right of the leftmost point of each slice.
For the smaller ordinate window in Theorem~\ref{thm:rectangle},
Section~\ref{sec:rational} proves positivity all the way from $4/5$
to $7/8$, so there is no such gap across that rectangle.
\end{remark}

\section{An exact rational rectangle certificate}
\label{sec:rational}

We now specialize to \eqref{eq:constants}, so $c=94/125$, and write
$S=S_E$. For $4/5\le\beta\le7/8$ define
\[
 p(\beta)=\left(\frac{219}{250}-\beta\right)
          \left(\frac{219}{250}+\beta-1\right),\qquad
 z=(\gamma-\Tcen)^2.
\]
The exact intervals needed are
\begin{equation}\label{eq:pz-box}
 p_-^{\rm box}=\frac{751}{1000000}\le p\le
 p_+^{\rm box}=\frac{3211}{62500},\qquad
 0\le z\le z_+^{\rm box}=\frac1{40000}.
\end{equation}
These bounds require no numerical monotonicity check. Indeed,
\begin{align*}
 p(\beta)-\frac{751}{1000000}
   &=\left(\frac78-\beta\right)\left(\beta-\frac18\right)\ge0,\\
 \frac{3211}{62500}-p(\beta)
   &=\left(\beta-\frac45\right)\left(\beta-\frac15\right)\ge0.
\end{align*}
The height bound $|\gamma-\Tcen|\le1/200$ gives the bound on $z$.

\begin{lemma}[Positivity on the entire parameter box]\label{lem:box}
For every real $(p,z)$ in \eqref{eq:pz-box}, $S(p,z)>0$.
\end{lemma}

\begin{proof}
The four values are the following positive rational numbers:
\begin{center}
\begin{tabular}{ccc}
\toprule
$p$ & $z$ & $S(p,z)$\\
\midrule
$751/1000000$ & $0$ & $1113147971/2000000000000$\\
$751/1000000$ & $1/40000$ & $185457099/500000000000$\\
$3211/62500$ & $0$ & $2828891/7812500000$\\
$3211/62500$ & $1/40000$ & $426282771/2000000000000$\\
\bottomrule
\end{tabular}
\end{center}
The quadratic part of $S$ is $-E(p-z)^2$. Its Hessian is
\[
 -2E\begin{pmatrix}1&-1\\-1&1\end{pmatrix},
\]
which is negative semidefinite. Thus $S$ is jointly concave.
Every point of the box is a convex combination of its four corners,
so concavity bounds $S$ below by the corresponding convex
combination of the corner values. In particular,
\begin{equation}\label{eq:margin}
 S(p,z)\ge\frac{426282771}{2000000000000}>0.
\end{equation}
The implemented Lean proof reaches strict positivity by concave
quadratic interpolation first in $p$ and then in $z$; it proves
positivity for all real points, not only for rational samples.
\end{proof}

\begin{proof}[Proof of Theorem~\ref{thm:rectangle}]
Suppose that a zero $s=\beta+i\gamma$ satisfied the stated rectangle
conditions. Since $\beta\ge4/5>0$, it cannot be a trivial zero.
Nontrivial zeros have $\beta<1$ and, by \eqref{eq:H},
$\beta\le7/8$. Equations \eqref{eq:pz-box} and
Lemma~\ref{lem:box} therefore place it in the region of
Theorem~\ref{thm:quartic}, which is impossible. The proof uses
strict positivity even at all four parameter-box corners, hence
the closed edges of the original rectangle are covered.
The identity $\zeta(\overline{s})=\overline{\zeta(s)}$ gives the
negative-height copy.
\end{proof}

\section{External numerical evaluation and reproducibility}
\label{sec:numeric}

Directly evaluating $\xi$ itself at the chosen ordinate is
unnecessary. Differentiating \eqref{eq:xi} gives the exact identity
\begin{equation}\label{eq:evaluation}
 \frac{\xi'}{\xi}(s)
 =\frac1s+\frac1{s-1}-\frac12\log\pi
  +\frac12\psi(s/2)+\frac{\zeta'(s)}{\zeta(s)},
\end{equation}
where $\psi=\Gamma'/\Gamma$. All displayed terms are retained in
the executable evaluation.

\begin{proposition}[External Arb enclosure]\label{prop:arb}
The supplied evaluation with Python~3.12.5, python-flint~0.9.0,
and linked FLINT~3.6.0
returns a real ball for the value $K$ in \eqref{eq:constants}
contained in the rational open interval
\begin{equation}\label{eq:Kinterval}
 \frac{1449534467251349699}{10^{17}}
 < K <
 \frac{1449534467251349700}{10^{17}}<\frac{29}{2}.
\end{equation}
This is a rigorous external arithmetic certificate, not a Lean
kernel theorem.
\end{proposition}

\begin{proof}[Computational verification]
Arb uses real and complex ball arithmetic with rigorous error
tracking \cite{JohanssonArb,PythonFlint}. The anchor is supplied
as the exact integer ordinate and a ball enclosing the exact
rational real part. A degree-one Taylor jet of $\zeta(s+X)$,
computed with rigorously bounded derivatives \cite{JohanssonZeta},
produces both $\zeta(s)$ and $\zeta'(s)$ in one calculation.
The program checks that the enclosure for $\zeta(s)$ does not
contain zero before dividing. Formula \eqref{eq:evaluation}
then gives a ball for $K$.

The evaluation is recomputed at 96 and 160 bits. Each run checks
strict interval separation from both rational endpoints in
\eqref{eq:Kinterval} and from $29/2$. At 160 bits the recorded
real ball is
\[
 \left[\,14.495344672513496991432147293120053955414637
          \ \pm\ 9.02\cdot10^{-43}\,\right].
\]
The full complex enclosures, precision settings, library version,
and successful rational comparisons are stored in the supplied
JSON record. The two precisions are replays of the same library;
they are not independent implementations of the special functions.
\end{proof}

The repository artifacts use the following paths:
\begin{center}
\begin{tabular}{>{\raggedright\arraybackslash}p{0.43\textwidth}>{\raggedright\arraybackslash}p{0.48\textwidth}}
\toprule
Artifact & Scope\\
\midrule
\src{scripts/local_shape_certificate.py}
 & Arb evaluation and exact rational geometry checks.\\[3pt]
\src{verification/local-shape-arb.json}
 & Recorded balls, exact constants, and corner values.\\[3pt]
\src{ZeroFree/LocalShape.lean}
 & Kernel-checked real algebra and conditional exclusions.\\
\bottomrule
\end{tabular}
\end{center}
With python-flint installed, reproduce the numerical calculation
from the repository root by running
\begin{center}
\texttt{python scripts/local\_shape\_certificate.py}
\end{center}
The script's optional plot is illustrative;
the certified geometry is the exact inequality and rational
rectangle, rather than a sampled contour or a decimal boundary.

\section{The formalization boundary}
\label{sec:formal}

The Lean module works with real variables and a generic predicate
$\mathsf{IsZero}(\beta,t)$. It does not define $\zeta$ or $\xi$.
Its analytic interface is an explicit hypothesis saying that a
putative zero obeys the paired kernel majorization with upper
bound $E$. In the mathematical application, Lemma~\ref{lem:pair}
and $K\le E$ supply this interface.

The module proves the exact four corner identities, positivity
throughout the real parameter box, the bounds taking a physical
rectangle into that box, the positive denominator, and the
identity \eqref{eq:pair-identity}. It then proves the generic
quartic exclusion, the rectangle exclusion, and the rectangle
exclusion with an explicit global-bound hypothesis. The
single-kernel open disk is also checked. The module additionally
proves finite nonnegative weighted budget inequalities, generic
conditional weighted exclusions, their reflected-pair specialization,
and the pointwise union limitation of Proposition~\ref{prop:limitation}.
Each analytic hypothesis
is a theorem parameter; no analytic claim is inserted as an
axiom declaration.

The axiom reports list only
\texttt{propext}, \texttt{Classical.choice}, and \texttt{Quot.sound}.
The module has no \src{sorry}, \src{sorryAx}, or
\src{native_decide}. These checks establish the finite algebra
and the conditional implications. They do not certify the
correctness of the external Arb library, formalize the
Hadamard product, or discharge \eqref{eq:H}. A future
end-to-end Lean theorem for the concrete zeta function must
connect and prove those separate inputs.

The written joint-concavity margin \eqref{eq:margin}, the
general symbolic parameters of Theorem~\ref{thm:quartic},
and the radical slice description are written analytic or
algebraic results of this paper. The current Lean module certifies
the fixed constants, strict real-box positivity, and abstract finite
weighted implications described above. In particular, the analytic
bridge from the concrete zeta function to the multi-anchor kernel
budgets remains external. The module does not claim to formalize
every generalized statement in the paper.

\section{Finite multi-anchor certificates and limitations}
\label{sec:multi}

The basic argument can use any finite collection of certified
anchors. This is useful when a common rational inequality is
easier to check on a whole box than a case split between
individual anchors.

\begin{theorem}[Nonnegative multi-anchor certificate]\label{thm:multi}
Assume that all nontrivial zeros satisfy $\beta\le B$, where
$1/2\le B<1$. For $1\le j\le m$ choose
$s_j=\sigma_j+iT_j$ with $\sigma_j>B$, and bounds
\[
 K_j=\Rea\frac{\xi'}{\xi}(s_j)\le E_j.
\]
Let $w_j\ge0$ be real weights, and set
\[
 \Phi_j(\sigma,t)=
 \frac{\sigma_j-\sigma}{(\sigma_j-\sigma)^2+(t-T_j)^2}
 +\frac{\sigma_j+\sigma-1}{(\sigma_j+\sigma-1)^2+(t-T_j)^2}.
\]
Then the region
\begin{equation}\label{eq:multi}
 \frac12<\sigma\le B,\qquad
 \sum_{j=1}^m w_j\Phi_j(\sigma,t)>
 \sum_{j=1}^m w_jE_j
\end{equation}
contains no zeros of $\zeta$.
\end{theorem}

\begin{proof}
For a hypothetical zero in the stated real-part range,
Lemma~\ref{lem:pair} at every anchor gives
$\Phi_j(\sigma,t)\le K_j\le E_j$. Multiplication by
nonnegative weights and finite summation contradict
\eqref{eq:multi}. All kernel denominators are positive
because $\sigma_j>B\ge\sigma>1/2$.
\end{proof}

\begin{proposition}[Pointwise limitation of nonnegative weights]
\label{prop:limitation}
For fixed anchors, bounds, and nonnegative weights, the region
\eqref{eq:multi} is contained in
\[
 \bigcup_{j=1}^m
 \left\{(\sigma,t):\frac12<\sigma\le B,
                         \ \Phi_j(\sigma,t)>E_j\right\}.
\]
\end{proposition}

\begin{proof}
If every $\Phi_j(\sigma,t)\le E_j$, their nonnegative
weighted sum cannot exceed the weighted sum of the $E_j$.
The contrapositive proves the assertion.
\end{proof}

Thus nonnegative weighting alone cannot create a pointwise
exclusion beyond the union of the individual certified
regions. Its role is to provide a common or compact
certificate for a box or a finite cover. A real improvement
beyond that union would require further information, such as
unavoidable contributions from other zeros or constraints on
a joint zero configuration. No such extra input is used here.

To extend the present geometry, one can first evaluate more
exact anchors and take the actual union of their reflected-pair
regions. Boxes inside this union can be certified by rational
polynomial inequalities and finite covering arguments. This
procedure gives local or finite-range statements; it does not
by itself prove a prescribed improvement at every ordinate.

There is also a separate analytic challenge in extracting
height-dependent improvements from the moment argument of
\cite{Cubic}. Its internal truncation height is not the
ordinate of a target zero, and a constant depending on a fixed
target cannot automatically be converted into an improvement
of an asymptotic exponent. A uniform curve such as
$B-c/\log|t|$ would require new quantitative input relating
the target ordinate to a strictly positive surplus in that
argument. The present shape certificate makes no such
inference.

Finally, a narrow high-ordinate rectangle can assist a
specific contour deformation or a finite-range estimate.
It does not lower a global error exponent in prime counting,
since the boundary outside the certified windows is unchanged.
The concrete contribution is the conditional quartic
exclusion region, its exact closed rectangle, the reproducible
external evaluation, and the finite real algebra verified in
Lean.

\begin{thebibliography}{9}

\bibitem{GoldsteinGrigutis}
E.~Gold\v{s}tein and A.~Grigutis,
\emph{On a positivity property of the real part of logarithmic
derivative of the Riemann $\xi$-function},
arXiv:2201.08599 (2022).
\url{https://arxiv.org/abs/2201.08599}.

\bibitem{JohanssonArb}
F.~Johansson,
\emph{Arb: efficient arbitrary-precision midpoint-radius interval
arithmetic}, IEEE Transactions on Computers \textbf{66}(8)
(2017), 1281--1292.
\url{https://doi.org/10.1109/TC.2017.2690633};
\url{https://arxiv.org/abs/1611.02831}.

\bibitem{PythonFlint}
The python-flint contributors,
\emph{python-flint 0.9.0 documentation},
real and complex ball arithmetic and power series documentation,
accessed October 11, 2026.
\url{https://python-flint.readthedocs.io/en/latest/}.

\bibitem{JohanssonZeta}
F.~Johansson,
\emph{Rigorous high-precision computation of the Hurwitz zeta
function and its derivatives}, arXiv:1309.2877 (2013).
\url{https://arxiv.org/abs/1309.2877}.

\bibitem{PlattTrudgian}
D.~Platt and T.~Trudgian,
\emph{The Riemann hypothesis is true up to $3\cdot10^{12}$},
arXiv:2004.09765 (2020).
\url{https://arxiv.org/abs/2004.09765}.

\bibitem{Cubic}
RH-Zero-Free-Formalization project,
\emph{A cubic zero-free boundary from plain-moment parameter feedback},
October 7, 2026, with the repository's explicit proof-status documentation.
\url{https://github.com/lixiang90/RH-Zero-Free-Formalization/blob/main/papers/kappa-feedback-cubic-boundary-paper.tex};
\url{https://github.com/lixiang90/RH-Zero-Free-Formalization/blob/main/docs/proof-status.md}.

\end{thebibliography}
\end{document}
